\section{Classical Sobolev embeddings and their realizations}\label{sec:critical-appendix}
The estimates below are consequences of classical Sobolev embedding.
For the Euclidean homogeneous inequality we use
Tao~\cite[Appendix~A, equation~(A.11)]{taodispersive}; for compact
manifolds we use Taylor~\cite[Chapter~13, Proposition~6.4 and the
preceding discussion]{taylor2011iii}. We retain the realization and
normalization details needed for our applications.
Throughout, $\Lambda=(-\Delta)^{1/2}$ has symbol $|\xi|$ on $\R^3$
and $2\pi|k|$ on the unit torus.

\begin{lemma}[Critical embeddings on the two domains]\label{lem:critical-embeddings}
For $a\in\{1/2,1\}$ and $p_a=6/(3-2a)$, there are finite constants,
depending only on $a$ and on the fixed unit torus when relevant, such that
\begin{equation}\label{eq:critical-embedding-pair}
 \norm{v}_{L^{p_a}(\R^3)}\le C_a\norm{\Lambda^a v}_{L^2(\R^3)},
 \qquad
 \norm{v}_{L^{p_a}(\TT)}\le C_a'\norm{\Lambda^a v}_{L^2(\TT)}.
\end{equation}
The first inequality holds initially for Schwartz functions and extends
to the homogeneous completion specified below.  The second holds for
mean-zero periodic distributions with finite displayed right-hand side;
the mean-zero condition is essential.  Componentwise application gives
the same statements for vector and tensor fields, with constants allowed
to depend on their fixed number of components.

Consequently, on either domain and whenever the indicated homogeneous
norms are finite (with mean-zero representatives on the torus),
\begin{align}
 \norm{v}_3&\le C\norm{v}_{\dot H^{1/2}},\nonumber\\
 \norm{\nabla v}_3+\norm{\Lambda v}_3
   &\le C\norm{v}_{\dot H^{3/2}},\label{eq:critical-derived}\\
 \norm{\nabla v}_6&\le C\norm{\Delta v}_2.\nonumber
\end{align}
Here the derivative statements on $\R^3$ refer to distributions whose
derivatives have the homogeneous realizations described in the proof;
in particular they apply to all the smooth $H^\infty$ fields used in
the article.
\end{lemma}

\begin{proof}
\emph{Euclidean embedding and the homogeneous realization.}
The classical homogeneous Sobolev inequality
\cite[Appendix~A, equation~(A.11)]{taodispersive}, with dimension three,
base exponent two, and $0<a<3/2$, gives
\[
 \norm{v}_{L^{p_a}(\R^3)}\le C_a\norm{\Lambda^a v}_2,
 \qquad \frac1{p_a}=\frac12-\frac a3.
\]
Tao's Fourier convention includes the compensating Plancherel factor
specified in Section~\ref{sec:intro}, so this is exactly the norm
used here. Initially take $\widehat v\in
C_c^\infty(\R^3\setminus\{0\})$ and complete this class in
$\norm{\Lambda^a v}_2$.

We specify its distributional realization to avoid any ambiguity by
polynomials. If $G\in L^2$ is the limit of $|\xi|^a\widehat v_n$, set
$\widehat v=|\xi|^{-a}G$. For every Schwartz test function $\phi$,
\[
 \int_{\R^3}|\xi|^{-a}|G(\xi)\widehat\phi(\xi)|\dd\xi
 \le\norm{G}_2
 \left(\int_{\R^3}|\xi|^{-2a}|\widehat\phi(\xi)|^2\dd\xi\right)^{1/2}
 <\infty.
\]
The last integral is finite at zero because $2a<3$, and at infinity
by rapid decay. This estimate also proves distributional convergence.
The embedding makes $v_n$ Cauchy in $L^{p_a}$, whose limit represents
the same distribution by H\"older's inequality. Every $G\in L^2$ is
approximable by smooth functions supported in compact annuli, by
truncation and mollification. Thus this construction gives a unique
$L^{p_a}$ representative for every element of the completion.
It includes all Schwartz functions and all $H^\infty$ fields with
finite displayed homogeneous norm. Taking $a=1/2$ and $a=1$ gives
$p_a=3$ and $p_a=6$, respectively.

\emph{The unit torus.}
The compact-manifold Sobolev embedding in
\cite[Chapter~13, Proposition~6.4 and the preceding discussion]{taylor2011iii}
yields
\[
 \norm{v}_{L^{p_a}(\TT)}\le C_a\norm{v}_{H^a(\TT)}
 \qquad(a=1/2,1).
\]
Taylor's angular torus coordinates are converted to the unit torus
as explained in Section~\ref{sec:intro}; the resulting inhomogeneous
norms are equivalent. For mean-zero fields the spectral gap gives
\[
 (1+4\pi^2|k|^2)^a
 \le(1+(2\pi)^{-2})^a(2\pi|k|)^{2a},\qquad k\ne0,
\]
and hence $\norm{v}_{H^a(\TT)}\le C_a'\norm{v}_{\dot H^a(\TT)}$.
This proves the periodic inequality. Fourier truncations converge in
$\dot H^a$ for every mean-zero periodic distribution with finite norm,
and therefore also in $L^{p_a}$; their distributional limits identify
the representatives.

\emph{Derivative estimates.}
Apply the $a=1/2$ inequality to each $\partial_jv$ and to $\Lambda v$,
and the $a=1$ inequality to each $\partial_jv$. Plancherel and
$\sum_j|\xi_j|^2=|\xi|^2$, or its Fourier-series counterpart, give
\eqref{eq:critical-derived}. The Euclidean statements use the preceding
realizations of the derivatives; no embedding of $\dot H^{3/2}$ into
$L^\infty$ is asserted. Componentwise application covers vectors and
tensors. In particular, for smooth fields,
\[
 |\ip{(v\cdot\nabla)v}{\Lambda v}|
 \le\norm{v}_3\norm{\nabla v}_3\norm{\Lambda v}_3
 \le C\norm{v}_{\dot H^{1/2}}\norm{v}_{\dot H^{3/2}}^2.
\]
All constants depend only on the fixed exponents, domain, and norm
conventions, not on the field or its frequency support.
\end{proof}
